\subsection{LIS3DH}
\label{chap:analysis_lis3dh}

Other elements intended for inclusion in the virtual laboratory, and which are 
part of the UPM laboratory equipment, are all the sensors from the Arduino Sensor 
Kit~\cite{seeedstudio_sensor_kit} by seeedStudio, including the LIS3DH accelerometer. This three-axis 
accelerometer, designed by STMicroelectronics, which has a digital output with an 
I2C interface, was selected in this thesis for integration into QEMU and to test the 
I2C interface of the emulated STM32F405 microcontroller.

\subsubsection{Characteristics}

The LIS3DH~\cite{LIS3DHDatasheet2016} is an ultra-low-power, high-performance, three-axis 
linear accelerometer belonging to the "nano" family, with a standard I2C/SPI digital 
serial interface output. The device features ultra-low-power operating modes for 
advanced energy savings and integrated intelligent functions, as well as a 32-level 
FIFO that can store up to 32 complete sets of samples.

The LIS3DH provides three-dimensional acceleration measurement with a 
selectable dynamic range of ±2 g, ±4 g, ±8 g, and ±16 g, depending on the required 
use case. The data output resolution is 16 bits. The device is capable of generating 
output data rates (ODR) from 1 Hz to 5.3 kHz, covering applications ranging from 
low-power real-time monitoring to high-frequency event capture. The current 
consumption of the LIS3DH varies significantly depending on the operating mode: in 
full operation mode, it can reach approximately 70 µA, while in low power modes, 
it drops to as low as 2 µA, a critical feature for battery-powered applications.

\subsubsection{Registers}

The LSI3DH has 31 eight-bit registers allocated to the address space. These 
registers are responsible for both device configuration and access to measurement 
data and status information. Each of these aspects is controlled by a different set 
of registers. This section of the chapter will analyze the most relevant ones. A table 
with the accelerometer's memory map can be found in Appendix \ref{ann:memories_maps}.

This STMelectronics accelerometer has three operating modes: low-power mode, 
normal mode, and high-resolution mode. These operating modes are defined by 
setting the LPen bit in CTRL\_REG1 and the HR bit in CTRL\_REG4, Table \ref{tab:lis3dh_modes}. 

\begin{table}[H]
    \centering
    \begin{tabular}{|l|c|c|l|l|}
        \hline
        \textbf{Mode} & \textbf{LPen} & \textbf{HR} & \textbf{Resolution} & \textbf{Characteristics} \\
        \hline
        Low-Power   &   1   & 0 &   8-bit   & Highest power efficiency,\\
                                            &&&&8-bit left-shifted data \\
        \hline
        Normal      &   0   & 0 &   10-bit  & Balanced power/performance,\\
                                            &&&&10-bit left-shifted data \\
        \hline
        High-Resolution & 0 & 1 &   12-bit  & Maximum precision,\\ 
                                            &&&& 12-bit left-shifted data \\
        \hline
    \end{tabular}
    \caption{LIS3DH operating modes and resolution}
    \label{tab:lis3dh_modes}
\end{table}

The output data rate is configured using the ODR[3:0] bits of the CTRL\_REG1 
register, table \ref{tab:lis3dh_odr}. This register also contains the individual enable bits (Xen, Yen, 
Zen) for each measurement axis.

\begin{table}[H]
    \centering
    \begin{tabular}{|cccc|l|}
        \hline
        \textbf{ODR3} & \textbf{ODR2} & \textbf{ODR1} & \textbf{ODR0} & \textbf{Power Mode / Output Data Rate} \\
        \hline
        0 & 0 & 0 & 0 & Power-down \\
        0 & 0 & 0 & 1 & HR/Normal/Low-power: 1~Hz \\
        0 & 0 & 1 & 0 & HR/Normal/Low-power: 10~Hz \\
        0 & 0 & 1 & 1 & HR/Normal/Low-power: 25~Hz \\
        0 & 1 & 0 & 0 & HR/Normal/Low-power: 50~Hz \\
        0 & 1 & 0 & 1 & HR/Normal/Low-power: 100~Hz \\
        0 & 1 & 1 & 0 & HR/Normal/Low-power: 200~Hz \\
        0 & 1 & 1 & 1 & HR/Normal/Low-power: 400~Hz \\
        1 & 0 & 0 & 0 & Low-power: 1.60~kHz \\
        1 & 0 & 0 & 1 & HR/Normal: 1.344~kHz; Low-power: 5.376~kHz \\
        \hline
    \end{tabular}
    \caption{lis3dh data rate configuration}
    \label{tab:lis3dh_odr}
\end{table}

As explained in its specifications, the LIS3DH supports four measurement ranges: 
±2 g, ±4 g, ±8 g, and ±16 g. These are configured using bits FS[1:0] in the 
CTRL\_REG4 register. The range affects both the maximum measurable acceleration 
and the sensitivity. Table \ref{tab:lis3dh_fs} summarizes this:

\begin{table}[H]
    \centering
    \begin{tabular}{|c|c|c|l|l|l|}
        \hline
        \textbf{FS[0:1]} & \textbf{Range} & \textbf{Range (mg)}
            & \textbf{$S_{0}$ Normal} & \textbf{$S_{0}$ High-Res} & \textbf{$S_{0}$ Low-Power} \\
        \hline
        00  &   ±2g     & ±2000 mg      & 4 mg/LSB      & 1 mg/LSB      &  16 mg/LSB\\
        01  &   ±4g     & ±4000 mg      & 8 mg/LSB      & 2 mg/LSB      &  32 mg/LSB\\
        10  &   ±8g     & ±8000 mg      & 16 mg/LSB     & 4 mg/LSB      &  64 mg/LSB\\
        11  &   ±16g    & ±16000 mg     & 48 mg/LSB     & 12 mg/LSB     &  192 mg/LSB\\
        \hline
    \end{tabular} 
    \caption{lis3dh full-scale ranges and sensitivities}
    \label{tab:lis3dh_fs}
\end{table}

The LSI3DH stores acceleration measurement data in the output data registers: 
OUT\_X\_L/H, OUT\_Y\_L/H, and OUT\_Z\_L/H. A very important feature of these 
registers is that their values are expressed in left-justified two's complement. Regarding 
the interruptions, the LIS3DH has two independent interrupt lines, INT1 and INT2. These can 
be configured to signal different types of events, significantly reducing the 
computational load on the host processor. The device supports three categories of 
interrupts: data interrupts (data ready and FIFO status), inertial event interrupts (axis 
movement, wake-up, freefall, and 6D/4D orientation detection), and click detection 
interrupts. The CTRL\_REG3 and CTRL\_REG6 control registers determine which 
events are routed to INT1 and INT2, respectively.